\subsection{势能方法的适用范围与稳定性}

势能函数也常记为 $U$。近地面重力势能 $mgy$、万有引力势能 $-GMm/r$、弹簧势能 $kx^2/2$
来自不同的力随位置的关系，不能跨越适用范围混用。尤其当高度与地球半径可比时，$g$ 随位置变化，$mgy$ 不再准确。

先选出已用势能 $U$ 表示的、与时间无关的保守相互作用；把所有尚未计入 $U$ 的力的功合记为 $W_{\rm other}$。
动能定理给出 $\Delta K=-\Delta U+W_{\rm other}$，所以
\begin{equation}
 \Delta(K+U)=W_{\rm other}.
 \label{ch08:eq:other-work}
\end{equation}
当 $W_{\rm other}=0$，机械能守恒。若还要把右边分为外力功和非保守内力功，两类必须互不重叠，
不能把同一个外部摩擦力计算两次。也不能既写重力功，又把同一重力势能变化列入左边。
在固定外场中可以把外场作用写入势能；例如研究近地面质点时使用 $mgy$，不必每次显式计算地球的运动。
但若场源也运动或相互作用显含时间，就须重新检查这个简化是否成立。

一维平衡点 $x_0$ 满足 $U'(x_0)=0$，即合力为零。若 $U''(x_0)>0$，势能在该点取局部极小，
微小偏离后出现指向平衡点的回复力，是稳定平衡；若 $U''(x_0)<0$，回复方向相反，平衡不稳定。
在二阶导数为正的平衡点附近展开：
\[
 U(x)\simeq U(x_0)+\frac12U''(x_0)(x-x_0)^2,
 \qquad F_x\simeq-U''(x_0)(x-x_0).
\]
与 $m\ddot x=-k_{\rm eff}(x-x_0)$ 比较可知，$k_{\rm eff}=U''(x_0)$，故小振动角频率
$\omega=\sqrt{U''(x_0)/m}$。如果 $U''(x_0)=0$，这个二阶判据不能定性，必须继续查看更高阶项；
例如 $U=kx^4/4$ 仍是稳定势阱，却不是上述线性回复模型。

给定一维总机械能 $E$，允许位置满足 $E\ge U(x)$；转折点满足 $E=U(x)$。在允许区间内
\[
 |v(x)|=\sqrt{\frac{2[E-U(x)]}{m}}.
\]
这只能先给出速率，运动方向由所处的往返分支确定。若需时间，在方向不变的一段上取
$v(x)=\pm\sqrt{2[E-U(x)]/m}$，再积分 $t_B-t_A=\int_{x_A}^{x_B}\dd x/v(x)$。
因此能量曲线既能判断可达区域，也能与运动方程配合求运动时间。
